% Generated by scripts/assemble.py; corrections belong in source/corrections.json.
\examhead{103}{浙江大学 2026 年数学分析试卷}{102}
% p103-b14; source PDF page 103
\par\Needspace{5\baselineskip}
\blockanchor{p103-b14}
\majorquestion{一}计算题. 需给出计算过程或依据. 共 40 分.\par
% p103-b15; source PDF page 103
\blockanchor{p103-b15}
\question{1}设 $x_{n} = \frac{1}{n + 1} +\frac{1}{n + 2} +\dots +\frac{1}{n + n}, n\in \mathbb{Z}^{+}$ ，求极限 $\lim_{n\to \infty}x_n.$\par
% p103-b16; source PDF page 103
\blockanchor{p103-b16}
\question{2}设 $f(x)$ 在区间 (0,2) 上可导, 且 $f(1) = f'(1) = 0$, $f''(1) = 2$, 求极限 $\lim_{x \to 1} \frac{xf'(x)}{xf'(x) - f(x)}$.\par
% p103-b17; source PDF page 103
\blockanchor{p103-b17}
\question{3}设 $x_{n} = \int_{0}^{\frac{\pi}{2}}\frac{\sin^{2}(nx)}{\sin x}\mathrm{d}x, n\in \mathbb{Z}^{+}$, 求极限 $\lim_{n\to \infty}\frac{x_n}{\ln n}$.\par
% p103-b18; source PDF page 103
\par\Needspace{6\baselineskip}
\blockanchor{p103-b18}
\question{4}设 $S$ 是 $x^{2} + y^{2} = z^{2}$ 介于 $z = 0$ 与 $z = 1$ 之间的部分, 并取外侧为正向, 计算第二型曲面积分\par
\nopagebreak[4]
% p103-b19; source PDF page 103
\blockanchor{p103-b19}
\displayformula{\iint_ {S} (x - 2 y + z) \mathrm{d} y \mathrm{d} z + (y - 2 z + x) \mathrm{d} z \mathrm{d} x + (z - 2 x + y) \mathrm{d} x \mathrm{d} y.}
% p103-b20; source PDF page 103
\par\Needspace{5\baselineskip}
\blockanchor{p103-b20}
\majorquestion{二}(15 分) 设 $f(x) = \begin{cases} x(1 - x^2), & x \in \mathbb{Q}, \\ x(1 + x^2), & x \notin \mathbb{Q}. \end{cases}$ 请讨论 $f(x)$ 在 $\mathbb{R}$ 上的连续性与可导性(需说明理由).\par
% p103-b21; source PDF page 103
\par\Needspace{5\baselineskip}
\blockanchor{p103-b21}
\majorquestion{三}(10 分) 证明: 设 $f(x)$ 在区间 $[0,1]$ 上只有第一类间断点, 则 $f(x)$ 是 $[0,1]$ 上的有界函数.\par
% p103-b22; source PDF page 103
\par\Needspace{5\baselineskip}
\blockanchor{p103-b22}
\majorquestion{四}(15 分) 设 $f(x)$ 在区间 $[0,1]$ 上二阶可导, 且 $f(0) = f(1) = 0$, 若 $\min_{x \in [0,1]} f(x) = -1$, 证明: 存在 $\xi \in (0,1)$, 使得 $f''(\xi) \geq 8$.\par
% p103-b23; source PDF page 103
\par\Needspace{5\baselineskip}
\blockanchor{p103-b23}
\majorquestion{五}(15 分) 证明: 函数 $y = \ln x$ 在 $[a, +\infty)$ 上一致连续, 其中 $a > 0$, 但在 $(0, +\infty)$ 上不一致连续.\par
% p104-b00; source PDF page 104
\par\Needspace{5\baselineskip}
\par\Needspace{6\baselineskip}
\blockanchor{p104-b00}
\majorquestion{六}(15 分) 设 $a_1 > 0, a_2 > 0$, 定义 $a_{n+2} = \sqrt{a_{n+1}} + \sqrt{a_n}, n = 1, 2, \cdots$, 证明: 数列 $\{a_n\}$ 收敛, 并求 $\lim_{n \to \infty} a_n$.
\majorquestion{七}(15 分) 设 $\lambda \geq 1$, 定义 $I(\lambda) = \int_0^1 \cos(\lambda x^2) \mathrm{d}x$, 证明: 存在与 $\lambda$ 无关的常数 $c > 0$, 使得 $|I(\lambda)| \leq \frac{c}{\sqrt{\lambda}}$.
\majorquestion{八}(15 分) 设 $f(x), g(x)$ 是区间 $(0,1)$ 上的有界函数, 定义\par
\nopagebreak[4]
% p104-b01; source PDF page 104
\blockanchor{p104-b01}
\displayformula{\varphi (x) = \sup _ {t \in (0, 1)} (f (t) + \sqrt {x} g (t)), x \in (0, 1).}
% p104-b02; source PDF page 104
\blockanchor{p104-b02}
证明: 存在常数 $C \geq 0$, 使得对任意 $x, y \in (0,1)$, 有 $|\varphi(x) - \varphi(y)| \leq C \sqrt{|x - y|}$.\par
% p104-b03; source PDF page 104
\par\Needspace{5\baselineskip}
\blockanchor{p104-b03}
\majorquestion{九}(10 分) 证明: 函数 $f(x) = \sum_{n=2}^{\infty} \frac{\sin nx}{n \ln n}$ 在 $x = 0$ 处连续.\par
